\section{The Season}

Jesus is speaking to his disciples about the last day. The sea is in tumult
and men are wasting away for fear of what is coming. Then they will see the
Son of Man coming in a cloud with power and majesty. Jesus tells his disciples
what to do when these things begin. “Look up and lift up your heads, because
your redemption is at hand” (Luke 21:28). That is the Gospel the Church reads
on the First Sunday of Advent.

\emph{Adventus} means coming or arrival. At Christmas we will adore the Son of
God born of Mary at Bethlehem. At the end of the age we will stand before him.
The face we will see then belongs to the one who has already come to find us
in our sins and give his life for us. We can ask him now to put right what our
sins have spoiled.

Israel had long prayed for God to come. Isaiah speaks to a people who have
known the shadow of death and are promised a great light. A child is born to
them and will sit upon the throne of David, establishing his kingdom in
judgment and justice (Isa. 9:2, 6--7). In another place the prophet pleads
that God would rend the heavens and come down (Isa. 64:1). Christians can pray
that plea because the kingdom has come near in Christ and we still pray for
its fullness. The Son took flesh from the very people to whom God first made
those promises. He was born in David's family.

At the Last Supper, as Jesus is about to leave his disciples, he promises them
another kind of coming. “If any one love me, he will keep my word. And my
Father will love him and we will come to him and will make our abode with him”
(John 14:23). He promises to be present in those who love him even after he
has gone to the Father.

St Bernard of Clairvaux calls this hidden coming the way from the first to the
last (Advent Sermon 5.1--3). In his first coming Christ was seen in the
weakness of our flesh. At his return all will see him in majesty. In the years
between, the one who redeemed us comes to dwell with those who keep his word.
Bernard asks what it means to keep that word. Bread on a shelf can go stale or
be nibbled by mice. Food we have eaten becomes part of us. Scripture is kept
in that way when it reaches the will and we begin to do what it says. Bernard
looks for the result in what the hand does, what the mouth says and what the
heart desires. The Lord's promise to make his dwelling concerns all three.
